% !TEX program = pdflatex
\documentclass[12pt]{article}

\usepackage[utf8]{inputenc}
\usepackage{CJKutf8}
\usepackage{amsmath,amssymb}
\usepackage{amsthm}
\usepackage{geometry}
\geometry{
  a4paper,
  top=25mm,
  bottom=25mm,
  left=20mm,
  right=20mm
}
\usepackage{xcolor}
\usepackage{url}
\usepackage[unicode,colorlinks=true,linkcolor=blue,urlcolor=blue]{hyperref}
\usepackage{xurl}
\Urlmuskip=0mu plus 2mu
\sloppy
\usepackage{tocloft}

\renewcommand{\cftsecleader}{\cftdotfill{\cftdotsep}}

\newtheorem{definition}{定义}[section]
\newtheorem{axiom}{公理}[section]
\newtheorem{assumption}{假设}[section]
\newtheorem{theorem}{定理}[section]
\newtheorem{proposition}{命题}[section]
\newtheorem{corollary}{推论}[section]
\newtheorem{lemma}{引理}[section]
\newtheorem{remark}{备注}[section]

\newcommand{\Jac}{\operatorname{Jac}}
\newcommand{\dd}{\mathrm{d}}
\newcommand{\code}[1]{\path{#1}}
\newcommand{\GFramework}{\textsc{G-Framework}}
\newcommand{\Ob}{\operatorname{Ob}}
\newcommand{\Cert}{\operatorname{Cert}}
\newcommand{\Loss}{\mathcal{L}}
\newcommand{\Flow}{\operatorname{Flow}}
\newcommand{\Upd}{\operatorname{Upd}}
\newcommand{\NN}{\mathbb{N}}
\newcommand{\ZZ}{\mathbb{Z}}
\newcommand{\RR}{\mathbb{R}}
\newcommand{\Hist}{\operatorname{Hist}}
\newcommand{\Crit}{\operatorname{Crit}}
\newcommand{\Enc}{\operatorname{Enc}}
\newcommand{\Code}{\operatorname{Code}}
\newcommand{\Src}{\mathrm{Src}}
\newcommand{\Ext}{\mathrm{Ext}}
\newcommand{\Unproved}{\mathrm{Unproved}}
\newcommand{\FalseTag}{\mathrm{False}}
\newcommand{\Verdict}{\operatorname{Verdict}}
\newcommand{\Car}{\operatorname{Car}}
\newcommand{\Img}{\operatorname{Im}}
\newcommand{\id}{\mathrm{id}}
\newcommand{\pr}{\operatorname{pr}}

\setlength{\emergencystretch}{6em}
\setlength{\parindent}{0pt}
\setlength{\parskip}{0.6em}

\begin{document}
\begin{CJK*}{UTF8}{gkai}

\title{基于扩展签发协议的主丛联络--路径积分--同伦修复塔\\与结构表示定理}
\newcommand{\CertificateAuthorCN}{Gao Zheng（高政）}
\author{\CertificateAuthorCN}
\date{2026-03-13}
\maketitle

\section*{前言}
\addcontentsline{toc}{section}{前言}

本文的目标不是孤立地给出一套扩展签发协议，而是在 \GFramework{} 的整体链条中，
把来源直签发、补证后签发、主丛联络、路径积分支撑类、同伦修复塔与结构表示定理
组织成一个可追踪的认证系统。若这些对象分散出现，它们容易只表现为若干独立的
证明技巧；本文将其压合为一条从命题来源、证书生成、障碍上升、修复递归到结构
编码的形式链条。

本文的第一层立场是签发语义的分层化。一个命题并不因为被写下就自动进入可用
知识库；它必须通过来源接口、补证接口或未证边界中的某一类判定。由此，签发不再
只是“承认结论”，而是将命题放入可审计的来源--证书--判定结构中。来源直签发给出
基础语义，补证后签发给出修复语义，未证保留给出边界语义，三者共同构成本文的
扩展签发协议。

本文的第二层立场是把“可用命题”与“可恢复证书”分开。来源命题若已有直接证书，
可以进入来源直签发；若只有目标陈述而缺少关键链条，则必须进入补证过程；若补证
仍不足，则只能保留为未证边界。这个分层避免把叙述性可信误写为证明性可信，也
避免把暂未补齐的命题直接并入后续定理调用。由此，本文提供的是一种证书治理机制：
每个结论不仅要有陈述，还要有来源、证书、判定值与可回放的推理位置。

本文的第三层立场是把修复塔写成认证动力学。主丛联络、路径积分与同伦扭曲在
本文中不是彼此松散的比喻，而是认证对象沿层级递归上升时的三种接口：联络给出
传输规则，路径积分给出支撑类聚合，同伦修复给出障碍吸收与正规形。于是，修复塔
不只是“补充证明”，而是一个逐层消化障碍、逐层生成证书、逐层恢复可签发性的
递归系统。

本文的第四层立场是主丛联络与路径积分的认证角色。主丛联络负责说明证书如何沿
对象层、性质层或路径层被传输；路径积分支撑类负责将多条可行路径的贡献压缩为
可比较的支撑对象；同伦扭曲则负责在遇到障碍时改变路径或结构的代表形。三者合在
一起，使签发协议不再只处理静态命题，而能够处理“命题在修复塔中如何被传输、
聚合、变形并重新获得证书”的动态过程。

本文的第五层立场是结构表示的可恢复性。单一实数总编码并不用于把数学对象压扁为
不可解释的数值，而是用于证明整座认证塔可以被无选择、可恢复、可追溯地表示。
因此，“一切皆结构”的弱版本在本文中被理解为：只要给定足够的生成元、投影、
证书链与恢复规则，集合元素就可以作为同调结构中的生成位置或投影位置被忠实追踪。

本文的第六层立场是把结构表示与签发协议闭合起来。若某个对象只能被编码而不能被
恢复，它就不能承担证书角色；若某个证书只能被局部验证而不能在塔中传递，它也不能
支撑扩展签发。本文因此同时要求编码、恢复、投影、证书与判定值相容。这样，结构
表示定理不只是附录式编码结果，而是整座扩展签发塔能够被机器化复验、跨层调用与
后续继承的基础条件。

本文的第七层立场是为后续同伦命中与 PDE 残差范式提供证书层接口。后续稿件中的
低残差命中、雅可比间隙下降、高阶正交方向与确定性极限，都需要知道哪些命题可以
作为已签发事实调用，哪些命题只能作为待补证边界保留。本文给出的扩展签发协议，
正是把“可调用事实”从叙述中剥离出来、放入证书系统的前置工作。

因此，本文在整体 \GFramework{} 中承担的是“签发--修复--结构表示”的认证锚定作用。
它向前连接可修复障碍与同伦类边界，向中间连接主丛联络、路径积分与递归修复塔，
向后为统计修复、同伦命中与 PDE 残差范式提供证书层接口。概括地说，本文把
“命题是否可用”这一问题，从口头可信性推进为可签发、可补证、可递归修复、可编码
恢复的形式系统。

更具体地说，本文给出的不是一份“证明已完成”的列表，而是一种管理证明状态的
协议。它允许命题处于来源直签发、补证后签发、未证保留等不同状态，并要求每一种
状态都有可追溯的判定值。这使得 \GFramework{} 后续文章在调用结论时，不必重新
猜测其证明强度，而可以直接沿签发塔找到证书来源。

从方法论上看，本文把“证书”从静态附件提升为动态对象。证书可以沿联络传输，可以
在路径积分支撑类中聚合，可以被障碍阻断，也可以通过同伦修复塔重新获得可签发性。
因此，扩展签发协议本身就是一种微型修复系统：它不只保存正确结论，也记录结论从
不可用到可用、从局部可证到全局可调用的路径。

从整体谱系看，本文处在基础修复与后续命中范式之间。前面的稿件说明障碍可以被
定义和修复；后面的稿件说明低残差集合可以被命中、轨迹可以被控制。本文则回答：
这些修复和命中过程中产生的命题，如何进入可复用知识库。这个位置使它成为
\GFramework{} 的证书流通层。

\setcounter{tocdepth}{2}
\setcounter{secnumdepth}{2}
\tableofcontents
\clearpage

\section{总述}

本文把此前讨论中得到的结果整理为一份自洽的定理化文稿。核心目标有四项：

其一，形式化 ``扩展签发'' 认证协议，并把来源直签发与补证后签发统一到同一框架之内；

其二，给出主丛联络、路径积分、障碍上升与同伦修复之间的递归塔式构造；

其三，证明整座塔存在显式、无选择、可恢复的单一实数总编码；

其四，证明集合元素可被忠实表示为同调结构中的生成元或投影，并据此给出 ``一切皆结构'' 的严格弱版本。

全文中，``严格'' 的含义是：每一结论都必须被写成对象明确、证书明确、推理链明确的数学命题；若某一结论尚缺关键基础证书，则只在备注中标出边界，而不把它并入定理结论。

\begin{remark}[阅读导航：仓内文稿与外部参照]
本文聚焦扩展签发协议、递归修复塔与结构表示定理。若读者需要单层“纤维终点空间--路径积分支撑类--同伦扭曲正规形”三重等价的前置形式化，可参照仓内笔记 \cite{RepoTex36Formal}；
若需要把“可修复障碍”作为对象层边界再向前追溯，可参照 \cite{RepoTex29RepairableObstruction}；
更系统的 law-space、主丛联络、路径积分与 Jacobiator-gap 背景，可参照本仓库卷册主稿 \cite{GaoZhengAppVol3}。

外部标准参照方面，主丛、联络、分类空间与联络存在性可参照 \cite{KobayashiNomizu1963Foundations,Husemoller1994FibreBundles}；
路径积分与半经典驻相写法可参照 \cite{FeynmanHibbs1965QMPI,Schulman2005Technique}；
同伦/同调与高阶障碍吸收的标准背景可参照 \cite{Hatcher2002AT,Whitehead1978Elements,Weibel1994HomologicalAlgebra,
  LodayVallette2012AlgebraicOperads}。
\end{remark}

\section{扩展签发认证协议}
\label{sec:tex35-protocol}

\subsection{命题空间与判定值}

\begin{definition}[命题空间与判定值]
记一切待判定命题构成集合
\[
\mathfrak P.
\]
对任意
\[
\Phi\in\mathfrak P,
\]
定义其判定值
\[
\Verdict(\Phi)\in\{\Src,\Ext,\Unproved\}.
\]
其中：
\[
\Src=\text{来源直签发},\qquad
\Ext=\text{扩展签发},\qquad
\Unproved=\text{未证明}.
\]
\end{definition}

\begin{definition}[来源直签发域]
定义
\[
\Sigma_{\mathrm{src}}\subseteq\mathfrak P
\]
为满足下述条件的命题集：存在只使用来源接口、来源定义、来源定理及类型一致直接推演链的证明
\[
\Pi_{\mathrm{src}}(\Phi)
\]
使得
\[
\Pi_{\mathrm{src}}(\Phi)\vdash \Phi.
\]
\end{definition}

\begin{definition}[扩展签发域]
定义
\[
\Sigma_{\mathrm{ext}}\subseteq\mathfrak P
\]
为满足下述条件的命题集：存在对象降型映射
\[
\mathfrak D_{\Phi},
\]
来源接口集
\[
\mathcal I_{\Phi},
\]
缺口集
\[
\mathcal G_{\Phi},
\]
新增证书包
\[
\mathcal C_{\Phi},
\]
以及补齐后的证明链
\[
\Pi_{\mathrm{ext}}(\Phi),
\]
使得
\[
(\mathfrak D_{\Phi},\mathcal I_{\Phi},\mathcal C_{\Phi},\Pi_{\mathrm{ext}}(\Phi))\vdash \Phi.
\]
\end{definition}

\begin{definition}[对象降型]
对象降型是把原句中的对象统一改写为同类型、可比较、可进证明链的对象。典型例子如下：
\[
P\rightsquigarrow \mathcal H_A(b_0,b_1),
\]
即把主丛总空间降为由联络 $A$ 筛出的可容许纤维演化类；
\[
\mathcal K_A\rightsquigarrow \Car(\mathcal K_A)=\mathcal P_A^{\mathrm{adm}},
\]
即把路径积分值降为其承载的可容许历史空间；
\[
\text{同伦扭曲}\rightsquigarrow \mathcal T_{A,H}^{\mathrm{eff}},
\]
即把模糊修辞降为有效同伦扭曲正规形证书类。
\end{definition}

\begin{theorem}[签发分层定理]
有
\[
\Sigma_{\mathrm{src}}\subseteq\Sigma_{\mathrm{ext}}.
\]
\end{theorem}

\begin{proof}
若 $\Phi\in\Sigma_{\mathrm{src}}$，则 $\Phi$ 已存在一条只用来源接口的证明链。此时取新增证书包为空：
\[
\mathcal C_{\Phi}=\varnothing.
\]
则 $\Phi$ 显然也满足扩展签发的定义条件，因此
\[
\Phi\in\Sigma_{\mathrm{ext}}.
\]
故
\[
\Sigma_{\mathrm{src}}\subseteq\Sigma_{\mathrm{ext}}.
\]
\end{proof}

\begin{proposition}[标准判定流程]
对任意新命题 $\Phi$，统一按如下流程判定：
\[
\Phi
\to
\mathfrak D_{\Phi}
\to
\mathcal I_{\Phi}
\to
\mathcal G_{\Phi}
\to
\mathcal C_{\Phi}
\to
\Pi(\Phi)
\to
\Verdict(\Phi).
\]
\end{proposition}

\begin{proof}
此流程只是把扩展签发定义展开。首先用对象降型消去类型不齐；随后抽取来源接口；然后定位剩余缺口；再以新增证书补齐；最后把完整证明链压缩成单一判定值。若任一步失败，则命题不能进入 $\Sigma_{\mathrm{ext}}$，只能标记为
  $\Unproved$。
\end{proof}

\section{来源接口的公理化摘录}

本节不引用对话叙述，而直接把后文所需的来源接口写成假设。它们不是新增结论，而是后续定理的输入。

\begin{assumption}[路径积分承载对象]\label{ass:path}
在 law-space 层上，存在可容许历史空间
\[
\Hist_n^{\mathrm{adm}}
\]
与作用量 $S_n$，使路径积分写作
\[
\mathcal K_n=
\int_{\Hist_n^{\mathrm{adm}}}
\exp\!\Bigl(\frac{i}{\hbar}S_n[\Gamma_n]\Bigr)\,\mathcal D\Gamma_n.
\]
并且 ordinary path integrals 可视为该 law-space 路径积分的投影。
这里保留的是经典路径积分的记号骨架及其“历史空间 + 作用量”表达，标准参照可见 \cite{FeynmanHibbs1965QMPI,Schulman2005Technique}；本文的 law-space 化重写与仓内递推语境可结合
  \cite{RepoTex36Formal,GaoZhengAppVol3} 一并理解。
\end{assumption}

\begin{assumption}[驻相历史]\label{ass:crit}
在半经典极限下，路径积分由驻相历史主导，即
\[
\Crit_n:=
\left\{[\Gamma_n]\in\Hist_n^{\mathrm{adm}}\;\middle|\;\frac{\delta S_n[\Gamma_n]}{\delta\Gamma_n}=0\right\}
\]
构成第 $n$ 阶路径积分的骨架。
\end{assumption}

\begin{assumption}[$n$ 阶 law-bundle]\label{ass:bundle}
对每一层级 $n\in\NN$，存在到第 $n$ 层可见的 law-bundle
\[
\pi_n:\mathcal E_n\to B_n,
\]
其水平提升尊重 $l_n$-可见不变量。
\end{assumption}

\begin{assumption}[总联络与统一 Bianchi]\label{ass:bianchi}
每一阶上允许几何层、参数层、法则层组成总联络
\[
\omega_{\mathrm{tot}}=\pr_{\mathrm{geom}}^*\omega_{\mathrm{geom}}+\pr_{\mathrm{par}}^*\omega_{\mathrm{par}}
  +\pr_{\mathrm{law}}^*\omega_{\mathrm{law}},
\]
其总曲率
\[
\mathcal F_{\mathrm{tot}}=\dd\omega_{\mathrm{tot}}+\frac12[\omega_{\mathrm{tot}}\wedge\omega_{\mathrm{tot}}]
\]
满足统一 Bianchi 恒等式
\[
D_{\mathcal A}\mathcal F_{\mathrm{tot}}=0.
\]
\end{assumption}

\begin{assumption}[障碍闭性与高阶修复]\label{ass:obstruction}
对每一阶 $n$ 与每一层级 $k$，若低阶括号
\[
l_1^{(n)},\dots,l_k^{(n)}
\]
已定义，则下一阶障碍
\[
\Ob_{k+1}^{(n)}
\]
满足
\[
\delta\Ob_{k+1}^{(n)}=0.
\]
若同调类
\[
[\Ob_{k+1}^{(n)}]=0,
\]
则存在某个边界型修复数据 $B_{k+1}^{(n)}$ 使
\[
\delta B_{k+1}^{(n)}=-\Ob_{k+1}^{(n)},
\]
从而可取
\[
l_{k+1}^{(n)}:=B_{k+1}^{(n)}
\]
完成高阶延拓与修复。
\end{assumption}

\begin{assumption}[法则层优先于表示层]\label{ass:prior}
原始 law-level 代数结构先于对象层 bundle 几何；对象层主丛解释与低维标量语义只是更高层 law-space 结构的表示或投影。
\end{assumption}

\section{后继封装、障碍上升与修复塔}
\label{sec:tex35-tower}

\subsection{基本定义}

\begin{definition}[后继基底]
在第 $n$ 阶上，取终点--语义等价关系 $\approx_n$，定义
\[
B_{n+1}:=\Crit_n/\!\approx_n.
\]
称 $B_{n+1}$ 为第 $n$ 阶驻相历史类的后继基底。
\end{definition}

\begin{definition}[后继联络]
若 $B_{n+1}$ 可赋予可接受的拓扑与光滑结构，并取结构群 $G_{n+1}$ 的分类映射
\[
c_{n+1}:B_{n+1}\to BG_{n+1},
\]
则定义主丛
\[
E_{n+1}:=c_{n+1}^*(EG_{n+1})\to B_{n+1}.
\]
若在该主丛上选定联络 $A_{n+1}$，则称 $(E_{n+1},A_{n+1})$ 为后继联络数据。
\end{definition}

\begin{definition}[第 $n$ 阶修复指数]
若对每个阶段 $n$ 存在最小的修复阶
\[
m_n:=\min\Bigl\{k\ge 1\;\Bigm|\;[\Ob_{k+1}^{(n)}]=0\ \text{且由}\ l_{k+1}^{(n)}\ \text{完成修复}\Bigr\},
\]
则称 $m_n$ 为第 $n$ 阶的修复指数。
\end{definition}

\subsection{三张证书}

\begin{proposition}[后继基底可定义证书 $C_{\mathrm{succ}}(n)$]\label{prop:succ}
若第 $n$ 阶驻相历史类 $\Crit_n$ 非空，且商空间
\[
B_{n+1}=\Crit_n/\!\approx_n
\]
可赋予 Hausdorff、paracompact、smooth 结构，则存在后继主丛
\[
E_{n+1}\to B_{n+1}
\]
与其联络 $A_{n+1}$。
\end{proposition}

\begin{proof}
商空间 $B_{n+1}$ 按定义存在。若它满足 Hausdorff、paracompact、smooth 条件，则可使用分类空间构造：选定分类映射
\[
c_{n+1}:B_{n+1}\to BG_{n+1},
\]
并令
\[
E_{n+1}:=c_{n+1}^*(EG_{n+1}).
\]
由于 paracompact 主丛上存在联络，故可取
\[
A_{n+1}\in\Omega^1(E_{n+1},\mathfrak g_{n+1}).
\]
这里用到的分类空间与主丛联络存在性，是标准纤维丛理论中的常见结论，可参照 \cite{KobayashiNomizu1963Foundations,Husemoller1994FibreBundles}。
于是后继主丛与后继联络同时存在。
\end{proof}

\begin{proposition}[障碍上升证书 $C_{\mathrm{obs}}(n,k)$]\label{prop:obs}
在第 $n$ 阶上，若前 $k$ 阶结构
\[
l_1^{(n)},\dots,l_k^{(n)}
\]
已定义，则下一阶残差形成障碍类
\[
\Ob_{k+1}^{(n)},
\]
并满足：
\[
\delta\Ob_{k+1}^{(n)}=0,
\qquad
[\Ob_{k+1}^{(n)}]=0\Longrightarrow \exists\, l_{k+1}^{(n)}\ \text{吸收该残差}.
\]
\end{proposition}

\begin{proof}
由假设 \ref{ass:obstruction}，障碍闭性是来源接口的一部分。其意义是：低阶 Stasheff/Jacobi 关系若只在截断层成立，则下一阶偏差自动成为 $\delta$-闭的候选障碍；这种“低阶失效
  $\Rightarrow$ 高阶障碍类 $\Rightarrow$ 边界吸收”的口径，与标准同伦代数/障碍论写法一致，可参照 \cite{Whitehead1978Elements,
  Weibel1994HomologicalAlgebra,LodayVallette2012AlgebraicOperads}。若该障碍类为零，则存在边界型修复数据 $B_{k+1}^{(n)}$ 使
\[
\delta B_{k+1}^{(n)}=-\Ob_{k+1}^{(n)}.
\]
定义
\[
l_{k+1}^{(n)}:=B_{k+1}^{(n)}
\]
即可把残差吸收到更高层结构中。这正是同伦修复塔的逐层上升规则。
\end{proof}

\begin{proposition}[逐阶编码证书 $C_{\mathrm{enc}}$]\label{prop:encn}
若每一阶存在修复指数 $m_n\in\NN$，则整座塔自动给出离散编码
\[
m=(m_n)_{n\in\NN}\in\NN^\NN.
\]
若再选定角变量
\[
\theta_n\in[0,2\pi),
\]
则逐阶定义
\[
r_n:=m_n+\frac{\theta_n}{2\pi}\in\RR,
\]
从而得到
\[
r=(r_n)_{n\in\NN}\in\RR^\NN.
\]
并且每一阶映射
\[
(m_n,\theta_n)\mapsto r_n
\]
为单射。
\end{proposition}

\begin{proof}
第一部分是直接的：每个阶段给出一个自然数修复指数 $m_n$，故得到一个自然数序列
\[
(m_n)_{n\in\NN}\in\NN^\NN.
\]
第二部分中，由
\[
r_n=m_n+\frac{\theta_n}{2\pi}
\]
可恢复
\[
m_n=\lfloor r_n\rfloor,
\qquad
\theta_n=2\pi(r_n-\lfloor r_n\rfloor).
\]
因此逐阶映射是单射。于是可同时得到 $\NN^\NN$ 码与 $\RR^\NN$ 码。
\end{proof}

\subsection{主定理}

\begin{theorem}[主丛联络--路径积分--同伦修复塔定理 $T_{\mathrm{tower}}$]\label{thm:tower}
设初始阶段给定
\[
(B_0,E_0\to B_0,A_0,S_0),
\]
并满足假设 \ref{ass:path}--\ref{ass:obstruction}。若对每个 $n\in\NN$：

\begin{itemize}
  \item 第 $n$ 阶可容许历史空间 $\Hist_n^{\mathrm{adm}}$ 可定义；
  \item 第 $n$ 阶驻相集 $\Crit_n$ 非空；
  \item 证书 $C_{\mathrm{succ}}(n)$ 与所有相关的 $C_{\mathrm{obs}}(n,k)$ 成立；
\end{itemize}
则存在一列递归构造的结构
\[
\mathfrak T=
\bigl(B_n,E_n\to B_n,A_n,\Hist_n^{\mathrm{adm}},\mathcal K_n,\{l_k^{(n)}\}_{k\ge 3}\bigr)_{n\in\NN}
\]
满足：
\[
\text{(i)}\quad B_{n+1}=\Crit_n/\!\approx_n,
\]
\[
\text{(ii)}\quad \mathcal K_n=
\int_{\Hist_n^{\mathrm{adm}}}\exp\!\Bigl(\frac{i}{\hbar}S_n[\Gamma_n]\Bigr)\,\mathcal D\Gamma_n,
\]
\[
\text{(iii)}\quad D_{A_n}\mathcal F_n=0,
\]
\[
\text{(iv)}\quad \text{任一低阶表观残差若障碍类为零，则可由某个更高 }l_{k+1}^{(n)}\text{ 吸收},
\]
从而 $\mathfrak T$ 构成以 $\NN$ 为离散指标的主丛联络--路径积分--同伦修复塔。
\end{theorem}

\begin{proof}
对 $n$ 作归纳。

基础步 $n=0$：由假设 \ref{ass:path}，第 $0$ 阶路径积分
\[
\mathcal K_0=
\int_{\Hist_0^{\mathrm{adm}}}\exp\!\Bigl(\frac{i}{\hbar}S_0[\Gamma_0]\Bigr)\,\mathcal D\Gamma_0
\]
存在；由假设 \ref{ass:crit}，驻相集 $\Crit_0$ 为其骨架。由命题 \ref{prop:succ} 得到
\[
B_1=\Crit_0/\!\approx_0,
\qquad
E_1\to B_1,
\qquad
A_1.
\]
若低阶截断出现残差，则由命题 \ref{prop:obs} 把它吸收到更高的 $l_{k+1}^{(0)}$ 中。因此第 $1$ 阶存在。

归纳步：设第 $n$ 阶已构成。由假设 \ref{ass:path} 与 \ref{ass:crit}，第 $n$ 阶的路径积分与驻相集存在。由命题 \ref{prop:succ} 构造
\[
B_{n+1},\ E_{n+1}\to B_{n+1},\ A_{n+1}.
\]
由假设 \ref{ass:bianchi}，总联络满足统一 Bianchi：
\[
D_{A_n}\mathcal F_n=0.
\]
若在低阶可见层上出现表观残差，则由命题 \ref{prop:obs} 上升到某个更高 $l_{k+1}^{(n)}$ 完成修复。于是第 $n+1$ 阶成立。

由数学归纳法，对一切 $n\in\NN$ 成立，定理得证。
\end{proof}

\begin{remark}[关于层级与维数]
定理 \ref{thm:tower} 中的 $n$ 是层级索引，不是流形维数。本文中 ``$n$ 阶'' 的含义是 ``到第 $n$ 层可见''，而不是 ``$n$ 维''。
\end{remark}

\section{单一实数总编码定理}
\label{sec:tex35-realcode}

本节承接主定理 \ref{thm:tower} 与逐阶编码证书 \ref{prop:encn}，把修复塔的逐层数据压缩为单一实数编码。

\begin{definition}[归一化映射]
定义
\[
\nu:\RR\to(0,1),\qquad
\nu(x)=\frac12+\frac{1}{\pi}\arctan x.
\]
则 $\nu$ 为双射，逆映射为
\[
\nu^{-1}(u)=\tan\bigl(\pi(u-\tfrac12)\bigr).
\]
\end{definition}

\begin{definition}[规范二进制展开]
对任意 $u\in(0,1)$，取其不以无限尾随 $1$ 结尾的二进制展开
\[
u=
\sum_{m=1}^{\infty} b_m(u)2^{-m},
\qquad
b_m(u)\in\{0,1\}.
\]
此约定保证展开唯一。
\end{definition}

\begin{definition}[Cantor 配对]
定义
\[
\langle n,m\rangle:=\frac12(n+m)(n+m+1)+m,
\qquad n,m\in\NN.
\]
它是 $\NN\times\NN\to\NN$ 的双射。
\end{definition}

\begin{theorem}[显式单一实数总编码定理]\label{thm:realcode}
存在一个显式、无选择、可恢复的单射
\[
\Enc_{\mathrm{tot}}:\RR^\NN\hookrightarrow \RR.
\]
更具体地，对任意
\[
x=(x_n)_{n\in\NN}\in\RR^\NN,
\]
令
\[
u_n:=\nu(x_n)=\sum_{m=1}^{\infty}b_{n,m}2^{-m},
\]
并定义
\[
\Enc_{\mathrm{tot}}(x):=
\sum_{n=0}^{\infty}\sum_{m=1}^{\infty}
 b_{n,m}\,2^{-(\langle n,m\rangle+1)}.
\]
则 $\Enc_{\mathrm{tot}}$ 为显式、无选择、可恢复的单一实数总编码。
\end{theorem}

\begin{proof}
先证定义良好。因为每个 $b_{n,m}\in\{0,1\}$，且
\[
\sum_{k\ge 1}2^{-k}\le 1,
\]
故双重级数绝对收敛，因此 $\Enc_{\mathrm{tot}}(x)\in[0,1]\subset \RR$。

再证单射。设
\[
\Enc_{\mathrm{tot}}(x)=\Enc_{\mathrm{tot}}(y).
\]
由二进制位的唯一性，在固定的规范展开约定下，所有编码位都相等：
\[
b_{n,m}(x)=b_{n,m}(y),\qquad \forall n,m.
\]
由 Cantor 配对函数的可逆性，可逐位恢复所有 $b_{n,m}$。于是对每个 $n$，有
\[
\nu(x_n)=\sum_{m\ge 1}b_{n,m}(x)2^{-m}=
\sum_{m\ge 1}b_{n,m}(y)2^{-m}=\nu(y_n).
\]
再由 $\nu$ 的可逆性得到
\[
x_n=y_n,\qquad \forall n.
\]
故 $x=y$。

最后证 ``无选择''。整个构造只使用显式公式 $\arctan$、规范二进制展开规则及 Cantor 配对函数，不调用选择公理，也不调用外部挑选函数。因此该编码是显式、无选择、可恢复的。
\end{proof}

\begin{corollary}[塔的单一实数总编码]\label{cor:towercode}
若修复塔 $\mathfrak T$ 已经由命题 \ref{prop:encn} 产生逐阶实数序列
\[
r(\mathfrak T)=(r_n)_{n\in\NN}\in\RR^\NN,
\]
则定义
\[
\Code(\mathfrak T):=\Enc_{\mathrm{tot}}(r(\mathfrak T))\in\RR.
\]
于是整座塔存在单一实数总编码。
\end{corollary}

\begin{proof}
直接把定理 \ref{thm:realcode} 应用于序列 $r(\mathfrak T)$ 即得。
\end{proof}

\begin{theorem}[拓扑自然编码不可能定理]\label{thm:notop}
若把 ``自然'' 解释为 ``保持原生拓扑结构的同胚编码''，则不存在
\[
\RR^\NN\cong \RR
\]
的拓扑自然总编码。
\end{theorem}

\begin{proof}
$\RR$ 是局部紧空间；而带乘积拓扑的 $\RR^\NN$ 不是局部紧空间。事实上，任取一点 $x=(x_n)$ 及其任意基本邻域
\[
U=\prod_{n\in\NN}U_n,
\]
其中除有限多个 $n$ 外均有 $U_n=\RR$。其闭包
\[
\overline U=\prod_{n\in\NN}\overline{U_n}
\]
仍含有无限多个非紧因子 $\RR$，故 $\overline U$ 不是紧的。因此 $\RR^\NN$ 的每个点都没有紧闭包邻域，所以 $\RR^\NN$ 不是局部紧空间。局部紧性在同胚下保持，故 $\RR^\NN$ 不可能与 $\RR$
  同胚。
\end{proof}

\begin{remark}[编码结论的边界]
定理 \ref{thm:realcode} 与推论 \ref{cor:towercode} 证明的是：存在显式、无选择、可恢复的单一实数总编码；定理 \ref{thm:notop} 同时说明：若把 ``自然'' 理解为
  ``保持原生拓扑的同胚''，则该要求不可满足。
\end{remark}

\section{结构表示定理：整数、同伦类与同调投影}
\label{sec:tex35-structure}

\begin{theorem}[整数 $1$ 的同伦--同调结构表示]\label{thm:one}
存在空间 $S^1$，使得
\[
[S^1,S^1]_*\cong \ZZ,
\qquad
H_1(S^1;\ZZ)\cong \ZZ,
\]
并且整数 $1$ 对应于基点保持同伦类中的度数为 $1$ 的元素，以及第一同调群中的标准生成元。
\end{theorem}

\begin{proof}
考虑覆盖映射
\[
p:\RR\to S^1,
\qquad
p(t)=e^{2\pi i t}.
\]
对任意基点保持映射 $f:(S^1,1)\to(S^1,1)$，存在唯一提升
\[
\widetilde f:\RR\to\RR,
\qquad
p\circ\widetilde f=f\circ p,
\qquad
\widetilde f(0)=0.
\]
由于 $p(t+1)=p(t)$，故
\[
p(\widetilde f(t+1))=p(\widetilde f(t)).
\]
于是
\[
\widetilde f(t+1)-\widetilde f(t)\in\ZZ.
\]
连续性迫使该差为常数，记作 $\deg(f)\in\ZZ$。由同伦提升性质可知，$\deg(f)$ 只依赖于同伦类。映射
\[
[f]_*\mapsto \deg(f)
\]
给出
\[
[S^1,S^1]_*\cong \ZZ.
\]
恒等映射 $\id_{S^1}$ 的提升是 $\widetilde f(t)=t$，因此其度数为 $1$。

再看同调。把 $S^1$ 看成一个含一个 $0$-胞腔和一个 $1$-胞腔的 CW 复形，则胞腔链群满足
\[
0\to C_1\cong \ZZ\xrightarrow{\partial_1}C_0\cong \ZZ\to 0,
\]
且两端附着到同一个 $0$-胞腔，故
\[
\partial_1=0.
\]
又无 $2$-胞腔，因此
\[
\operatorname{im}\partial_2=0.
\]
从而
\[
H_1(S^1;\ZZ)=\ker\partial_1/\operatorname{im}\partial_2\cong \ZZ.
\]
其标准 $1$-胞腔类即为生成元，对应整数 $1$。
\end{proof}

\begin{theorem}[任意集合元素的同调结构表示]\label{thm:allstructure}
对任意集合 $A$，存在拓扑空间
\[
X_A:=\bigvee_{a\in A}S_a^1,
\]
使得
\[
H_1(X_A;\ZZ)\cong \bigoplus_{a\in A}\ZZ.
\]
并且对每个 $a\in A$，同时存在：
\[
j_A(a)=e_a\in H_1(X_A;\ZZ),
\]
即对应的坐标生成元，以及
\[
\pi_a:H_1(X_A;\ZZ)\to \ZZ,
\qquad
\pi_a\Bigl(\sum_{b\in A}n_be_b\Bigr)=n_a,
\]
即对应的坐标投影。映射
\[
j_A:A\hookrightarrow H_1(X_A;\ZZ)
\]
与
\[
\iota_A:A\hookrightarrow \operatorname{Hom}(H_1(X_A;\ZZ),\ZZ),
\qquad
\iota_A(a)=\pi_a,
\]
均为单射。
\end{theorem}

\begin{proof}
构造
\[
X_A=\bigvee_{a\in A}S_a^1,
\]
即以一个公共基点把以 $A$ 为指标的一族圆粘接而成的花束空间。它有一个 $0$-胞腔和对每个 $a\in A$ 的一个 $1$-胞腔 $e_a$，因此胞腔链群为
\[
C_1(X_A;\ZZ)\cong \bigoplus_{a\in A}\ZZ e_a,
\qquad
C_0(X_A;\ZZ)\cong \ZZ.
\]
因为每个 $1$-胞腔的两个端点都附着到同一个公共基点，所以
\[
\partial_1(e_a)=0,
\qquad \forall a\in A.
\]
又由于没有 $2$-胞腔，故
\[
\operatorname{im}\partial_2=0.
\]
于是
\[
H_1(X_A;\ZZ)=\ker\partial_1/\operatorname{im}\partial_2\cong \bigoplus_{a\in A}\ZZ e_a.
\]
定义
\[
j_A(a):=e_a.
\]
若 $a\neq b$，则 $e_a\neq e_b$，故 $j_A$ 为单射。

再定义坐标投影
\[
\pi_a\Bigl(\sum_{b\in A}n_be_b\Bigr)=n_a.
\]
若 $a\neq b$，则
\[
\pi_a(e_a)=1,\qquad \pi_b(e_a)=0,
\]
所以 $\pi_a\neq\pi_b$，从而
\[
\iota_A(a):=\pi_a
\]
也是单射。故任意集合元素都可被忠实表示为某个同调结构中的生成元或投影。
\end{proof}

\begin{corollary}[``一切皆结构'' 的严格弱版本]\label{cor:weakstructure}
任意集合 $A$ 的任意元素 $a\in A$，都不必作为空壳原子处理；它总可以被忠实表示为某个富结构中的可区分对象。更精确地说，存在某个结构对象 $S_A$ 及单射
\[
\iota_A:A\hookrightarrow \mathcal O(S_A)
\]
使得 $a$ 被表示为 $\mathcal O(S_A)$ 中的生成元、投影或其他结构性可区分对象。
\end{corollary}

\begin{proof}
由定理 \ref{thm:allstructure}，可取
\[
S_A:=H_1(X_A;\ZZ),
\]
并令 $\mathcal O(S_A)$ 取为其生成元族或投影族。于是集合元素被忠实嵌入到结构对象中。故 ``一切皆结构'' 在 ``忠实结构表示'' 的意义下成立。
\end{proof}

\begin{remark}[边界]
推论 \ref{cor:weakstructure} 证明的是 ``元素可被结构忠实表示''，而不是 ``元素与某一唯一结构绝对同一''。后者需要额外的唯一性与自然性证书，不是本文结论的一部分。
\end{remark}

\section{相对结构基础定理与边界}
\label{sec:tex35-foundation}

本节承接定理 \ref{thm:allstructure} 与推论 \ref{cor:weakstructure}，讨论该结构表述何时可以上升为基础层结论，以及其严格边界。

\begin{theorem}[相对结构基础定理]\label{thm:relativefoundation}
在假设 \ref{ass:prior} 下，主丛联络--路径积分--同伦修复塔至少可作为来源框架内部的相对结构基础。更具体地：法则层结构是原像，对象层主丛几何与低维标量语义是其表示或投影，因此塔式结构在来源框架内部具有统摄性的上层地位。
\end{theorem}

\begin{proof}
由假设 \ref{ass:prior}，原始 law-level 结构先于对象层 bundle 几何；因此对象层主丛只是在更高法则空间中的表示。又由同一假设，低维标量语义只是更高 law-space 结构的扁平投影，而非第一性对象。

于是，在来源框架内部，具有逻辑先行地位的是法则层与修复塔式结构，而非对象层几何或标量语义。换言之，bundle 几何与 scalar semantics 都由更高结构统一生成、统一约束、统一解释。因此该塔至少可作为来源框架内部的相对结构基础。

\end{proof}

\begin{definition}[全局基础替代所需证书]
若某理论 $T$ 要被认证为 ``已全局替代集合论作为全部基础''，则至少需要给出如下证书包：
\[
F_1:\ \text{内部宇宙或对象域构造};
\]
\[
F_2:\ \text{元素关系或等价的类型归属/子对象机制};
\]
\[
F_3:\ \text{空、对、并、函数、幂对象或等价闭包机制};
\]
\[
F_4:\ \text{自然数对象与无穷机制};
\]
\[
F_5:\ \text{分离、替换、归纳或等价构造原则};
\]
\[
F_6:\ \text{将通常数学解释进 }T\text{ 的解释定理或保守性定理}.
\]
\end{definition}

\begin{proposition}[全局基础替代的必要条件]\label{prop:foundationneed}
若理论 $T$ 尚未给出上述 $F_1$--$F_6$，则不能把 ``$T$ 已全局替代集合论作为全部基础'' 认证为已证明结论。
\end{proposition}

\begin{proof}
所谓 ``全局替代集合论作为全部基础''，并不是只要求结构丰富，而是要求 $T$ 足以承载 ordinary mathematics 的全部基础功能。若缺少内部宇宙、子对象机制、基本闭包、自然数对象、替换或解释定理中的任意关键一项，
  就无法证明普通数学能被完全解释进 $T$。因此在缺少 $F_1$--$F_6$ 的情况下，不能把该结论视为已证明。
\end{proof}

\begin{corollary}[当前边界]
在本文范围内，可以严格签发的结论是：
\[
\text{该塔在来源框架内部构成相对结构基础。}
\]
而下述更强命题仍不并入正文定理结论：
\[
\text{该塔已全局替代集合论或 ZFC 作为全部基础。}
\]
\end{corollary}

\begin{proof}
前半句由定理 \ref{thm:relativefoundation} 给出；后半句由命题 \ref{prop:foundationneed} 限定其边界。
\end{proof}

\section{结论}
\label{sec:tex35-conclusion}

本文完成了以下几项严格结果。

第一，在第 \ref{sec:tex35-protocol} 节建立了扩展签发认证协议。它把 ``来源直签发'' 与 ``新增证书后签发'' 统一到同一判定机制中，并给出对象降型、来源接口、缺口集、证书包与证明链的标准流程。

第二，在第 \ref{sec:tex35-tower} 节基于假设 \ref{ass:path}、\ref{ass:bianchi} 与 \ref{ass:obstruction} 构造了以 $\NN$
  为离散指标的主丛联络--路径积分--同伦修复塔，并由定理 \ref{thm:tower} 给出主结论。

第三，在第 \ref{sec:tex35-realcode} 节由定理 \ref{thm:realcode} 与推论 \ref{cor:towercode} 证明了该塔存在显式、无选择、可恢复的单一实数总编码；同时由定理 \ref{thm:notop} 说明：若把 ``自然'' 理解为 ``保持原生拓扑的同胚''，则此类编码不可能存在。

第四，在第 \ref{sec:tex35-structure} 节由定理 \ref{thm:one}、定理 \ref{thm:allstructure} 与推论 \ref{cor:weakstructure} 证明了整数 $1$
  可以由度数为 $1$ 的同伦类或第一同调群中的标准生成元表示；更一般地，任意集合元素都可被忠实表示为某个同调结构中的生成元或投影，从而得到 ``一切皆结构'' 的严格弱版本。

第五，在第 \ref{sec:tex35-foundation} 节由定理 \ref{thm:relativefoundation} 证明了该塔在来源框架内部构成相对结构基础；同时由命题 \ref{prop:foundationneed}
  明确指出，若要把它升级为全局基础替代，还需要补齐内部宇宙、子对象机制、基本闭包、自然数对象、替换原理及解释定理等基础证书。

\section*{参考文献}
\begin{thebibliography}{99}

\bibitem{RepoTex36Formal}
【仓内 TeX】主丛联络的纤维终点空间、路径积分支撑类与同伦扭曲正规形三重等价的形式系统笔记：
\code{NoteBook/notebook2/tex36_pub/gframework_principal_connection_fiber_space_path_integral_homotopy_twist_equivalence_
  formal_system.tex}。

\bibitem{RepoTex29RepairableObstruction}
【仓内 TeX】可修复障碍同伦类与广义分形的形式系统笔记：
\code{NoteBook/notebook2/tex29_pub/gframework_repairable_obstruction_homotopy_class_vs_generalized_fractal_formal_system
  .tex}。

\bibitem{GaoZhengAppVol3}
Gao, Z. (2025).
\newblock GaoZheng G-Framework and GaoZheng G-Algebra: Applied Mathematics Volume III – HACA, PACER, and
  Certificate-Based AI.
\newblock In \emph{GaoZheng G-Framework and GaoZheng G-Algebra: Applied Mathematics Volume III – HACA, PACER, and
  Certificate-Based AI (v1.0)}.
\newblock Zenodo.
\newblock \url{https://doi.org/10.5281/zenodo.17762295}.


\bibitem{KobayashiNomizu1963Foundations}
Kobayashi, S., \& Nomizu, K. (1963).
\newblock \emph{Foundations of Differential Geometry, Volume I}.
\newblock Interscience Publishers.

\bibitem{Husemoller1994FibreBundles}
Husemoller, D. (1994).
\newblock \emph{Fibre Bundles} (3rd ed.).
\newblock Springer.

\bibitem{FeynmanHibbs1965QMPI}
Feynman, R. P., \& Hibbs, A. R. (1965).
\newblock \emph{Quantum Mechanics and Path Integrals}.
\newblock McGraw-Hill.

\bibitem{Schulman2005Technique}
Schulman, L. S. (2005).
\newblock \emph{Techniques and Applications of Path Integration}.
\newblock Dover Publications.

\bibitem{Hatcher2002AT}
Hatcher, A. (2002).
\newblock \emph{Algebraic Topology}.
\newblock Cambridge University Press.

\bibitem{Whitehead1978Elements}
Whitehead, G. W. (1978).
\newblock \emph{Elements of Homotopy Theory}.
\newblock Springer.

\bibitem{Weibel1994HomologicalAlgebra}
Weibel, C. A. (1994).
\newblock \emph{An Introduction to Homological Algebra}.
\newblock Cambridge University Press.

\bibitem{LodayVallette2012AlgebraicOperads}
Loday, J.-L., \& Vallette, B. (2012).
\newblock \emph{Algebraic Operads}.
\newblock Springer.
\end{thebibliography}

\end{CJK*}
\end{document}
